\section*{Overview}

Strongly connected components are the natural building blocks of directed graphs with cycles. Inside one SCC, every
vertex can reach every other. Once those components are collapsed, the graph becomes a DAG, and many problems become
much easier.

\section*{When to Use It}

Use SCC decomposition when:

\begin{itemize}
  \item the graph is directed and reachability inside cycles matters,
  \item the problem asks about mutual reachability or cycle groups,
  \item a graph with cycles should be reduced to a DAG before DP or counting.
\end{itemize}

\section*{Core Idea}

Kosaraju's algorithm uses two DFS passes:

\begin{itemize}
  \item one DFS on the original graph to record finishing order,
  \item one DFS on the reversed graph in reverse finishing order to extract components.
\end{itemize}

\section*{Key Insight}

The finishing order from the first pass places SCCs in a useful reverse-topological arrangement. When the second pass
starts from the latest unfinished vertex on the reversed graph, it captures exactly one SCC at a time.

\section*{Operations / Main Technique}

\begin{itemize}
  \item build the reverse graph,
  \item first DFS for finishing order,
  \item second DFS for component assignment,
  \item optionally build the condensation DAG.
\end{itemize}

\paragraph{Worked example.}
If vertices \(\{1,2,3\}\) can all reach each other, they form one SCC no matter how many internal edges they have. A
single edge leaving that group becomes one edge in the condensation DAG.

\section*{Correctness Intuition}

In the condensation DAG, finishing times in the first pass place source components late enough that the reversed-graph
DFS can collect them without leaking into another unassigned component. Repeating that process partitions the graph
exactly into SCCs.

\section*{Complexity Analysis}

Kosaraju runs in \(O(n + m)\) time and \(O(n + m)\) memory.

\section*{Implementation}

The code returns:

\begin{itemize}
  \item a component id for each vertex,
  \item the list of vertices in each component.
\end{itemize}

That is enough to build a condensation graph afterward if needed.

\section*{Common Pitfalls}

\begin{itemize}
  \item Forgetting to build the reversed graph.
  \item Reusing the visited array incorrectly between the two DFS passes.
  \item Confusing SCCs with weakly connected components.
\end{itemize}

\section*{Variants / Extensions}

\begin{itemize}
  \item Tarjan's one-pass SCC algorithm.
  \item 2-SAT by SCC on the implication graph.
  \item DAG DP after condensation.
\end{itemize}

\section*{Practice Problems}

\begin{itemize}
  \item Count SCCs in a directed graph.
  \item Build and analyze the condensation DAG.
  \item Solve 2-SAT or reachability constraints through SCC decomposition.
\end{itemize}

\section*{References}

\begin{itemize}
  \item \href{https://cp-algorithms.com/graph/strongly-connected-components.html}{cp-algorithms: Strongly Connected Components and Condensation Graph}.
  \item \href{https://usaco.guide/adv/SCC?lang=cpp}{USACO Guide: Strongly Connected Components}.
  \item \href{https://codeforces.com/catalog}{Codeforces Catalog}.
\end{itemize}
